\chapter{The RTL-to-GDS ASIC Flow}
\label{ch:flow}

This chapter gives the overall methodology --- the sequence of design
steps that turn a Verilog description into a manufacturable layout --- and
the tools and technology used at each stage. The remaining chapters then
apply this flow to the 1-bit slice and to the two 32-bit designs.

\section{Flow overview}
The RTL-to-GDS (or ``RTL-to-GDSII'') flow is the standard digital ASIC
implementation methodology. Design intent is captured at the
register-transfer level and progressively refined --- through synthesis
and physical implementation --- into geometry, while functional and
physical correctness are checked at every hand-off. \Cref{fig:flow}
summarises the steps used in this project and the artefacts produced at
each one.

\begin{figure}[htbp]
\centering
\begin{tikzpicture}[node distance=6mm and 22mm,font=\small]
  \node[flow] (spec) {1.\ Specification\\{\footnotesize function table, I/O}};
  \node[flow,below=of spec] (rtl) {2.\ RTL design\\{\footnotesize Verilog: structural + behavioural}};
  \node[flow,below=of rtl] (sim) {3.\ Functional simulation\\{\footnotesize Xcelium + testbenches}};
  \node[flow,below=of sim] (syn) {4.\ Logic synthesis\\{\footnotesize Genus + \texttt{gsclib045}}};
  \node[flow,below=of syn] (pnr) {5.\ Place \& Route\\{\footnotesize Innovus: floorplan, place, route}};
  \node[flow,below=of pnr] (sign) {6.\ Physical signoff\\{\footnotesize DRC, connectivity, timing}};
  \node[flow,below=of sign] (gds) {7.\ GDSII export\\{\footnotesize core layout stream-out}};
  \node[flow,below=of gds] (chip) {8.\ Chip integration\\{\footnotesize pad frame in Virtuoso}};

  \foreach \a/\b in {spec/rtl,rtl/sim,sim/syn,syn/pnr,pnr/sign,sign/gds,gds/chip}
     \draw[flowarr] (\a) -- (\b);

  % Right-hand annotations (artefacts)
  \node[right=of rtl,align=left,font=\footnotesize] (a1)
       {\code{.v} source \& testbenches};
  \node[right=of syn,align=left,font=\footnotesize] (a2)
       {gate netlist, \code{.sdc},\\ area/power/timing reports};
  \node[right=of pnr,align=left,font=\footnotesize] (a3)
       {routed \code{DEF}, parasitics};
  \node[right=of gds,align=left,font=\footnotesize] (a4)
       {\code{alu\_32bit.gds}};
  \foreach \n/\a in {rtl/a1,syn/a2,pnr/a3,gds/a4}
     \draw[-{Latex[length=1.6mm]},gray,dashed] (\n.east) -- (\a.west);

  % Feedback loop: signoff failures return to place & route / synthesis
  \draw[flowarr,accent!70,dashed]
     (sign.west) -- ++(-1.2,0) |- node[pos=0.25,left,font=\footnotesize,align=right]
        {ECO / fix\\violations} (pnr.west);
\end{tikzpicture}
\caption{The RTL-to-GDS flow applied in this project. Solid arrows show
the forward flow; the dashed feedback path indicates iteration when
signoff checks (timing, DRC, connectivity) fail.}
\label{fig:flow}
\end{figure}

\section{Process technology and libraries}
All implementation uses the Cadence \textbf{GPDK045} generic
\SI{45}{\nano\meter} CMOS process design kit. \Cref{tab:tech} lists the
key technology and library settings that appear throughout the synthesis
and PnR reports.

\begin{table}[htbp]
\centering
\caption{Technology and library configuration.}
\label{tab:tech}
\small
\begin{tabularx}{\linewidth}{@{}l X@{}}
\toprule
Item & Setting \\
\midrule
Process               & Cadence GPDK045, \SI{45}{\nano\meter} bulk CMOS \\
Standard-cell library & \texttt{gsclib045} (LEF: \texttt{gsclib045\_tech.lef},
                        \texttt{gsclib045\_macro.lef}) \\
Timing corner         & \texttt{PVT\_0P9V\_125C}, worst-case
                        \texttt{slow\_vdd1v0} liberty \\
Supply / ground nets  & \texttt{VDD} / \texttt{VSS} \\
Routing pin layer     & Metal~4 (from the pin-assignment file) \\
Multi-mode multi-corner& \texttt{mmmc.view} (setup on slow corner) \\
\bottomrule
\end{tabularx}
\end{table}

Choosing the \emph{slow}, low-voltage, high-temperature corner
(\SI{0.9}{\volt}, \SI{125}{\celsius}) for setup analysis is deliberate: it
represents the worst case for cell delay, so timing that closes here
closes across the operating range. All area/power/timing numbers quoted in
this report are from this corner unless stated otherwise.

\section{Tools used}
\Cref{tab:tools} maps each flow step to the Cadence tool that performs it.

\begin{table}[htbp]
\centering
\caption{Tools used at each stage of the flow.}
\label{tab:tools}
\small
\begin{tabularx}{\linewidth}{@{}l l X@{}}
\toprule
Stage & Tool & Role \\
\midrule
Simulation      & Xcelium (SimVision) & Event-driven verification of RTL against testbenches \\
Synthesis       & Genus 21.17         & RTL $\rightarrow$ gate netlist mapped to \texttt{gsclib045}; area/power/timing \\
Place \& Route  & Innovus             & Floorplan, power, placement, routing, post-route timing \\
Signoff         & Innovus / Pegasus / Assura & DRC and connectivity (LVS-style) checks \\
Chip assembly   & Virtuoso            & Pad-frame placement, I/O wiring, full-chip stream-out \\
\bottomrule
\end{tabularx}
\end{table}

\section{Design flavours implemented}
Because the assignment requires the 32-bit ALU to be built ``in two
ways,'' three distinct RTL designs are carried through the flow:
\begin{description}[leftmargin=2.4em,style=nextline]
  \item[Part 1 --- 1-bit ALU] the bit-slice cell, verified in isolation
        and taken through synthesis and PnR on its own
        (\cref{ch:part1}).
  \item[Part 2a --- 32-bit hierarchical] a \emph{structural} design that
        instantiates thirty-two 1-bit slices and connects their carries in
        a ripple chain (\cref{ch:part2a}).
  \item[Part 2b --- 32-bit behavioural] a single-module design that models
        the operations directly on 32-bit vectors with a \code{casez}
        statement (\cref{ch:part2b}).
\end{description}
The two 32-bit cores are compared in \cref{ch:compare}; the smaller and
faster one is streamed out and dropped into the pad frame in
\cref{ch:pad}.
